\chapter{The Arithmetic of M. de Mazarin.}

\lettrine{W}{hilst} the king was directing his course rapidly towards  the wing of the castle occupied by the cardinal, taking nobody with him but his  valet de chambre, the officer of musketeers came out, breathing like a man who  has for a long time been forced to hold his breath, from the little cabinet of  which we have already spoken, and which the king believed to be quite solitary.  This little cabinet had formerly been part of the chamber, from which it was  only separated by a thin partition. It resulted that this partition, which was  only for the eye, permitted the ear the least indiscreet to hear every word  spoken in the chamber.

There was no doubt, then, that this lieutenant of musketeers had heard all that  passed in his majesty's apartment.

Warned by the last words of the young king, he came out just in time to salute  him on his passage, and to follow him with his eyes till he had disappeared in  the corridor.

Then as soon as he had disappeared, he shook his head after a fashion  peculiarly his own, and in a voice which forty years' absence from  Gascony had not deprived of its Gascon accent, <A melancholy  service,> said he, <and a melancholy master!>

These words pronounced, the lieutenant resumed his place in his fauteuil,  stretched his legs and closed his eyes, like a man who either sleeps or  meditates.

During this short monologue and the mise en scene that had accompanied it,  whilst the king, through the long corridors of the old castle, proceeded to the  apartment of M. de Mazarin, a scene of another sort was being enacted in those  apartments.

Mazarin was in bed, suffering a little from the gout. But as he was a man of  order, who utilized even pain, he forced his wakefulness to be the humble  servant of his labour. He had consequently ordered Bernouin, his valet de  chambre, to bring him a little travelling-desk, so that he might write in bed.  But the gout is not an adversary that allows itself to be conquered so easily;  therefore, at each movement he made, the pain from dull became sharp.

<Is Brienne there?> asked he of Bernouin.

<No, monseigneur,> replied the valet de chambre; <M. de  Brienne, with your permission, is gone to bed. But if it is the wish of your  eminence, he can speedily be called.>

<No, it is not worth while. Let us see, however. Cursed ciphers!>

And the cardinal began to think, counting on his fingers the while.

<Oh, ciphers is it?> said Bernouin. <Very well! if your  eminence attempts calculations, I will promise you a pretty headache to-morrow!  And with that please to remember M. Guenaud is not here.>

<You are right, Bernouin. You must take Brienne's place, my friend.  Indeed, I ought to have brought M. Colbert with me. That young man goes on very  well, Bernouin, very well; a very orderly youth.>

<I do not know,> said the valet de chambre, <but I  don't like the countenance of your young man who goes on so well.>

<Well, well, Bernouin! We don't stand in need of your advice. Place  yourself there: take the pen and write.>

<I am ready, monseigneur; what am I to write?>

<There, that's the place: after the two lines already  traced.>

<I am there.>

<Write seven hundred and sixty thousand livres.>

<That is written.>

<Upon Lyons\longdash> The cardinal appeared to hesitate.

<Upon Lyons,> repeated Bernouin.

<Three millions nine hundred thousand livres.>

<Well, monseigneur?>

<Upon Bordeaux, seven millions.>

<Seven?> repeated Bernouin.

<Yes,> said the cardinal, pettishly, <seven.> Then,  recollecting himself, <You understand, Bernouin,> added he,  <that all this money is to be spent?>

<Eh! monseigneur; whether it be spent or put away is of very little  consequence to me, since none of these millions are mine.>

<These millions are the king's; it is the king's money I am  reckoning. Well, what were we saying? You always interrupt me!>

<Seven millions upon Bordeaux.>

<Ah! yes; that's right. Upon Madrid four millions. I give you to  understand plainly to whom this money belongs, Bernouin, seeing that everybody  has the stupidity to believe me rich in millions. I repel the silly idea. A  minister, besides, has nothing of his own. Come, go on. Rentrées generales,  seven millions; properties, nine millions. Have you written that,  Bernouin?>

<Yes, monseigneur.>

<Bourse, six hundred thousand livres; various property, two millions. Ah!  I forgot—the furniture of the different châteaux\longdash>

<Must I put of the crown?> asked Bernouin.

<No, no; it is of no use doing that—that is understood. Have you  written that, Bernouin?>

<Yes, monseigneur.>

<And the ciphers?>

<Stand straight under one another.>

<Cast them up, Bernouin.>

<Thirty-nine millions two hundred and sixty thousand livres,  monseigneur.>

<Ah!> cried the cardinal, in a tone of vexation; <there are  not yet forty millions!>

Bernouin recommenced the addition.

<No, monseigneur; there want seven hundred and forty thousand  livres.>

Mazarin asked for the account, and revised it carefully.

<Yes, but,> said Bernouin, <thirty-nine millions two hundred  and sixty thousand livres make a good round sum.>

<Ah, Bernouin; I wish the king had it.>

<Your eminence told me that this money was his majesty's.>

<Doubtless, as clear, as transparent as possible. These thirty-nine  millions are bespoken, and much more.>

Bernouin smiled after his own fashion—that is, like a man who believes no  more than he is willing to believe—whilst preparing the cardinal's  night draught, and putting his pillow to rights.

<Oh!> said Mazarin, when the valet had gone out; <not yet  forty millions! I must, however, attain that sum, which I had set down for  myself. But who knows whether I shall have time? I sink, I am going, I shall  never reach it! And yet, who knows that I may not find two or three millions in  the pockets of my good friends the Spaniards? They discovered Peru, those  people did, and—what the devil! they must have something left.>

As he was speaking thus, entirely occupied with his ciphers, and thinking no  more of his gout, repelled by a preoccupation which, with the cardinal, was the  most powerful of all preoccupations, Bernouin rushed into the chamber, quite in  a fright.

<Well!> asked the cardinal, <what is the matter now?>

<The king, monseigneur,—the king!>

<How?—the king!> said Mazarin, quickly concealing his paper.  <The king here! the king at this hour! I thought he was in bed long ago.  What is the matter, then?>

The king could hear these last words, and see the terrified gesture of the  cardinal rising up in his bed, for he entered the chamber at that moment.

<It is nothing, monsieur le cardinal, or at least nothing which can alarm  you. It is an important communication which I wish to make to your eminence  to-night,—that is all.>

Mazarin immediately thought of that marked attention which the king had given  to his words concerning Mademoiselle de Mancini, and the communication appeared  to him probably to refer to this source. He recovered his serenity then  instantly, and assumed his most agreeable air, a change of countenance which  inspired the king with the greatest joy; and when Louis was seated,—

<Sire,> said the cardinal, <I ought certainly to listen to  your majesty standing, but the violence of my complaint\longdash>

<No ceremony between us, my dear monsieur le cardinal,> said Louis  kindly: <I am your pupil, and not the king, you know very well, and this  evening in particular, as I come to you as a petitioner, as a solicitor, and  one very humble, and desirous to be kindly received, too.>

Mazarin, seeing the heightened colour of the king, was confirmed in his first  idea; that is to say, that love thoughts were hidden under all these fine  words. This time, political cunning, as keen as it was, made a mistake; this  colour was not caused by the bashfulness of a juvenile passion, but only by the  painful contraction of the royal pride.

Like a good uncle, Mazarin felt disposed to facilitate the confidence.

<Speak, sire,> said he, <and since your majesty is willing  for an instant to forget that I am your subject, and call me your master and  instructor, I promise your majesty my most devoted and tender  consideration.>

<Thanks, monsieur le cardinal,> answered the king; <that  which I have to ask of your eminence has but little to do with myself.>

<So much the worse!> replied the cardinal; <so much the  worse! Sire, I should wish your majesty to ask of me something of importance,  even a sacrifice; but whatever it may be that you ask me, I am ready to set  your heart at rest by granting it, my dear sire.>

<Well, this is what brings me here,> said the king, with a beating  of the heart that had no equal except the beating of the heart of the minister;  <I have just received a visit from my brother, the king of  England.>

Mazarin bounded in his bed as if he had been put in relation with a Leyden jar  or a voltaic pile, at the same time that a surprise, or rather a manifest  disappointment, inflamed his features with such a blaze of anger, that Louis  XIV, little diplomatist as he was, saw that the minister had hoped to hear  something else.

<Charles II.?> exclaimed Mazarin, with a hoarse voice and a  disdainful movement of his lips. <You have received a visit from Charles  II.?>

<From King Charles II.,> replied Louis, according in a marked  manner to the grandson of Henry IV the title which Mazarin had forgotten to  give him. <Yes, monsieur le cardinal, that unhappy prince has touched my  heart with the relation of his misfortunes. His distress is great, monsieur le  cardinal, and it has appeared painful to me, who have seen my own throne  disputed, who have been forced in times of commotion to quit my  capital,—to me, in short, who am acquainted with misfortune,—to  leave a deposed and fugitive brother without assistance.>

<Eh!> said the cardinal, sharply; <why had he not, as you  have, a Jules Mazarin by his side? His crown would then have remained  intact.>

<I know all that my house owes to your eminence,> replied the king,  haughtily, <and you may well believe that I, on my part, shall never  forget it. It is precisely because my brother, the king of England has not  about him the powerful genius who has saved me, it is for that, I say, that I  wish to conciliate the aid of that same genius, and beg you to extend your arm  over his head, well assured, monsieur le cardinal, that your hand, by touching  him only, would know how to replace upon his brow the crown which fell at the  foot of his father's scaffold.>

<Sire,> replied Mazarin, <I thank you for your good opinion  with regard to myself, but we have nothing to do yonder: they are a set of  madmen who deny God, and cut off the heads of their kings. They are dangerous,  observe, sire, and filthy to the touch after having wallowed in royal blood and  covenantal murder. That policy has never suited me,—I scorn it and reject  it.>

<Therefore you ought to assist in establishing a better.>

<What is that?>

<The restoration of Charles II, for example.>

<Good heavens!> cried Mazarin, <does the poor prince flatter  himself with that chimera?>

<Yes, he does,> replied the young king, terrified at the  difficulties opposed to this project, which he fancied he could perceive in the  infallible eye of his minister; <he only asks for a million to carry out  his purpose.>

<Is that all—a little million, if you please!> said the  cardinal, ironically, with an effort to conquer his Italian accent. <A  little million, if you please, brother! Bah! a family of mendicants!>

<Cardinal,> said Louis, raising his head, <that family of  mendicants is a branch of my family.>

<Are you rich enough to give millions to other people, sire? Have you  millions to throw away?>

<Oh!> replied Louis XIV, with great pain, which he, however, by a  strong effort, prevented from appearing on his countenance;—<oh!  yes, monsieur le cardinal, I am well aware I am poor, and yet the crown of  France is worth a million, and to perform a good action I would pledge my crown  if it were necessary. I could find Jews who would be willing to lend me a  million.>

<So, sire, you say you want a million?> said Mazarin.

<Yes, monsieur, I say so.>

<You are mistaken, greatly mistaken, sire; you want much more than  that,—Bernouin!—you shall see, sire, how much you really  want.>

<What, cardinal!> said the king, <are you going to consult a  lackey about my affairs?>

<Bernouin!> cried the cardinal again, without appearing to remark  the humiliation of the young prince. <Come here, Bernouin, and tell me  the figures I gave you just now.>

<Cardinal, cardinal! did you not hear me?> said Louis, turning pale  with anger.

<Do not be angry, sire; I deal openly with the affairs of your majesty.  Every one in France knows that; my books are as open as day. What did I tell  you to do just now, Bernouin?>

<Your eminence commanded me to cast up an account.>

<You did it, did you not?>

<Yes, my lord.>

<To verify the amount of which his majesty, at this moment, stands in  need. Did I not tell you so? Be frank, my friend.>

<Your eminence said so.>

<Well, what sum did I say I wanted?>

<Forty-five millions, I think.>

<And what sum could we find, after collecting all our resources?>

<Thirty-nine millions two hundred and sixty thousand.>

<That is correct, Bernouin; that is all I wanted to know. Leave us  now,> said the cardinal, fixing his brilliant eye upon the young king,  who sat mute with stupefaction.

<However\longdash> stammered the king.

<What, do you still doubt, sire?> said the cardinal. <Well,  here is a proof of what I said.>

And Mazarin drew from under his bolster the paper covered with figures, which  he presented to the king, who turned away his eyes, his vexation was so deep.

<Therefore, as it is a million you want, sire, and that million is not  set down here, it is forty-six millions your majesty stands in need of. Well, I  don't think that any Jews in the world would lend such a sum, even upon  the crown of France.>

The king, clenching his hands beneath his ruffles, pushed away his chair.

<So it must be then!> said he; <my brother the king of  England will die of hunger.>

<Sire,> replied Mazarin, in the same tone, <remember this  proverb, which I give you as the expression of the soundest policy:  <Rejoice at being poor when your neighbour is poor likewise.>>

Louis meditated this for a few moments, with an inquisitive glance directed to  the paper, one end of which remained under the bolster.

<Then,> said he, <it is impossible to comply with my demand  for money, my lord cardinal, is it?>

<Absolutely, sire.>

<Remember, this will secure me a future enemy, if he succeed in  recovering his crown without my assistance.>

<If your majesty only fears that, you may be quite at ease,>  replied Mazarin, eagerly.

<Very well, I say no more about it,> exclaimed Louis XIV

<Have I at least convinced you, sire?> placing his hand upon that  of the young king.

<Perfectly.>

<If there be anything else, ask it, sire; I shall most happy to grant it  to you, having refused this.>

<Anything else, my lord?>

<Why yes; am I not devoted body and soul to your majesty? Hola!  Bernouin!—lights and guards for his majesty! His majesty is returning to  his own chamber.>

<Not yet, monsieur: since you place your good-will at my disposal, I will  take advantage of it.>

<For yourself, sire?> asked the cardinal, hoping that his niece was  at length about to be named.

<No, monsieur, not for myself,> replied Louis, <but still for  my brother Charles.>

The brow of Mazarin again became clouded, and he grumbled a few words that the  king could not catch.  